\documentclass[12pt,a4paper]{article}
\usepackage[utf8]{inputenc}
\usepackage{geometry}
\geometry{margin=1in}
\usepackage[english]{babel}
\usepackage{setspace}
\onehalfspacing
\usepackage{titlesec}
\usepackage{graphicx}
\usepackage{booktabs}
\usepackage{longtable}
\usepackage{amsmath,amssymb}
\usepackage{hyperref}
\usepackage{xcolor}
\usepackage{float}
\usepackage{caption}
\usepackage{subcaption}
\usepackage{enumitem}
\usepackage{listings}
\usepackage{fancyhdr}
\usepackage[numbers]{natbib}

\hypersetup{
    colorlinks=true,
    linkcolor=blue!70!black,
    citecolor=blue!70!black,
    urlcolor=blue!70!black
}

\lstset{
    language=Python,
    basicstyle=\small\ttfamily,
    keywordstyle=\color{blue},
    stringstyle=\color{red!60!black},
    commentstyle=\color{green!60!black},
    numbers=left,
    numberstyle=\tiny\color{gray},
    breaklines=true,
    frame=single,
    backgroundcolor=\color{gray!5}
}

\pagestyle{fancy}
\fancyhf{}
\rhead{Market Price Forecasting for Ugandan Staple Crops}
\lhead{AI \& ML Project Report}
\cfoot{\thepage}

\begin{document}

% ============================================================
% TITLE PAGE
% ============================================================
\begin{titlepage}
\centering
\vspace*{1.5cm}

{\Large\textsc{Uganda Christian University}}\\[0.3cm]
{\large Faculty of Engineering, Design and Technology}\\[0.2cm]
{\large Department of Computer Science}\\[1cm]

\rule{\textwidth}{1.5pt}\\[0.4cm]
{\LARGE\bfseries Market Price Forecasting for Ugandan Staple Crops}\\[0.2cm]
{\large A Multi-Model Time-Series Approach Using WFP Market Data and Climate Covariates}\\[0.2cm]
\rule{\textwidth}{1.5pt}\\[1.5cm]

{\large\textbf{Project Report}}\\[0.3cm]
{\large Artificial Intelligence and Machine Learning}\\[0.3cm]
{\large Master of Science in Computer Science (Year 1, Semester 2)}\\[2cm]

\begin{tabular}{lll}
\textbf{Name} & \textbf{Reg. Number} & \textbf{Student Number} \\
\midrule
Jamugisa Peter Paul & S25M25/008 & B35503 \\
ILEGO Cynthia Cecilia & S25M25/007 & B35502 \\
AYEBARE Moses & S25M25/003 & B35498 \\
Apilli Barbra Ebal & S25M25/001 & B35496 \\
\end{tabular}

\vfill
{\large April 2025}
\end{titlepage}

% ============================================================
% ABSTRACT
% ============================================================
\newpage
\begin{abstract}
Price volatility for staple crops -- maize, beans, and cassava -- poses a significant threat to food security and economic stability among smallholder farmers in Uganda. This project develops and compares three time-series forecasting models for predicting crop market prices 4 to 8 weeks ahead at district level. We train ARIMA (SARIMAX), Facebook Prophet, and Long Short-Term Memory (LSTM) neural networks on historical price data from the World Food Programme (WFP), augmented with CHIRPS satellite rainfall data and World Bank commodity indices as exogenous covariates. Feature engineering includes price lags, rolling statistics, sinusoidal seasonal encoding, and climate lag features. Models are evaluated on three crops (Maize, Beans, Cassava flour) across three districts (Kampala, Gulu, Mbale) using Mean Absolute Error (MAE), Root Mean Squared Error (RMSE), and Mean Absolute Percentage Error (MAPE) on a chronological 80/20 test split. The LSTM model achieves the best average performance (MAE: 503.13 UGX, RMSE: 606.40, MAPE: 17.65\%), outperforming ARIMA (MAPE: 25.11\%) and Prophet (MAPE: 29.66\%) by substantial margins. The best-performing model is deployed in an interactive Streamlit web dashboard enabling stakeholders to generate real-time price forecasts. This work demonstrates that multivariate deep learning approaches, when combined with domain-relevant climate data, can meaningfully improve agricultural price prediction in data-scarce East African contexts.

\noindent\textbf{Keywords:} time-series forecasting, LSTM, ARIMA, Prophet, crop price prediction, food security, Uganda, machine learning, Streamlit
\end{abstract}

% ============================================================
% TABLE OF CONTENTS
% ============================================================
\newpage
\tableofcontents
\newpage

% ============================================================
% 1. INTRODUCTION
% ============================================================
\section{Introduction}

\subsection{Background and Motivation}
Agriculture contributes approximately 24\% of Uganda's Gross Domestic Product (GDP) and employs over 70\% of the national workforce \cite{worldbank2023}. Staple crops -- maize, beans, and cassava -- form the backbone of both food security and rural incomes. However, market prices for these commodities exhibit substantial volatility driven by weather variability, seasonal harvest cycles, global commodity shocks, and infrastructure limitations.

Price volatility creates a ``poverty trap'' for smallholder farmers: unable to predict future prices, farmers either sell immediately post-harvest when prices are lowest or hold produce hoping for price increases, risking spoilage. Early price signals -- reliable forecasts 4 to 8 weeks ahead -- would enable farmers to optimise selling decisions, allow NGOs to pre-position food aid before price spikes, and support evidence-based policy planning by the Ministry of Agriculture, Animal Industry and Fisheries (MAAIF).

Despite the importance of agricultural markets, existing monitoring systems by the World Food Programme (WFP) and Famine Early Warning Systems Network (FEWS NET) are largely retrospective, reporting current and recent prices rather than generating predictive forecasts at district-level granularity. This project addresses that analytical gap.

\subsection{Problem Statement}
No reliable, open-source, data-driven forecasting tool currently provides district-level price projections for Ugandan staple crops. This project builds and compares three machine learning forecasting models -- ARIMA, Facebook Prophet, and LSTM -- to determine which approach achieves the best predictive performance when trained on multivariate data including historical prices, satellite rainfall, and global commodity indices.

\subsection{Research Question}
\textit{How accurately can machine learning models forecast staple crop market prices 4 to 8 weeks ahead in Uganda, and which model -- ARIMA, Prophet, or LSTM -- achieves the best performance when rainfall is incorporated as a predictor?}

\subsection{Objectives}
\begin{enumerate}
    \item Collect, clean, and merge WFP historical price data with CHIRPS rainfall and World Bank commodity price indices for selected Ugandan districts.
    \item Engineer temporal, seasonal, and climate-based features as model inputs.
    \item Implement and train three forecasting models: ARIMA (SARIMAX), Facebook Prophet, and LSTM.
    \item Evaluate and compare model performance using MAE, RMSE, and MAPE on a held-out test set.
    \item Deploy the best-performing model in an interactive web dashboard built with Streamlit.
\end{enumerate}

\subsection{Scope}
The project covers three crops (Maize, Beans, Cassava flour) across three districts (Kampala, Gulu, Mbale) using data from 2010 to 2024. Kampala represents the central urban market, Gulu the Northern Corridor post-conflict region, and Mbale the Eastern agricultural belt with cross-border trade dynamics.

% ============================================================
% 2. LITERATURE REVIEW
% ============================================================
\section{Literature Review}

\subsection{Time-Series Forecasting in Agricultural Markets}
Agricultural price forecasting has evolved from simple moving average models to sophisticated machine learning approaches. \citet{zhang2003} demonstrated that hybrid ARIMA-ANN models outperform either approach alone for financial time series predicting. \citet{xiong2015} applied deep learning to agricultural commodity prices, showing that LSTM networks capture nonlinear dynamics missed by linear models.

\subsection{ARIMA Models}
The Autoregressive Integrated Moving Average (ARIMA) model, formalised by \citet{box2015}, remains the gold standard for univariate time-series forecasting. ARIMA models assume linear relationships between past values, past errors, and future values. The model is specified by three parameters $(p, d, q)$: the autoregressive order $p$, the degree of differencing $d$, and the moving average order $q$. Extensions include SARIMAX, which incorporates seasonal components and exogenous variables.

\subsection{Facebook Prophet}
Prophet, developed by \citet{taylor2018}, decomposes time series into trend, seasonality, and holiday components using an additive model. It is designed for business forecasting at scale and handles missing data and outliers gracefully. Prophet uses a piecewise linear trend with automatic changepoint detection and Fourier series for seasonal patterns.

\subsection{LSTM Neural Networks}
Long Short-Term Memory networks, introduced by \citet{hochreiter1997}, are a class of recurrent neural networks (RNNs) designed to learn long-range dependencies. Unlike standard RNNs that suffer from vanishing gradients, LSTMs use gated mechanisms (forget, input, and output gates) to selectively retain and update information across time steps. Recent applications in crop price forecasting include \citet{sabu2020}, who demonstrated superior LSTM performance on Indian commodity markets.

\subsection{Climate Covariates in Price Modelling}
Uganda's agriculture is predominantly rain-fed, making climate data a theoretically strong predictor of agricultural output and, consequently, market prices. \citet{brown2009} established the link between CHIRPS rainfall anomalies and food security outcomes in East Africa. Our work extends this by incorporating rainfall as an explicit exogenous predictor in the forecasting pipeline.

\subsection{Research Gap}
While time-series forecasting for commodity prices is well-studied globally, no published work provides a systematic comparison of ARIMA, Prophet, and LSTM at district-level granularity for East African staple crops using multivariate data that includes satellite rainfall. This project fills that gap.

% ============================================================
% 3. METHODOLOGY
% ============================================================
\section{Methodology}

\subsection{Research Design}
This project follows a quantitative experimental design. Three forecasting models are trained on identical datasets and evaluated using the same metrics and evaluation protocol, enabling direct and fair comparison.

\subsection{Data Sources}

\subsubsection{WFP Food Prices}
The primary dataset was obtained from the World Food Programme's Vulnerability Analysis and Mapping (VAM) data repository via the Humanitarian Data Exchange (HDX). It contains market-level price observations for multiple commodities across Ugandan districts, recorded at weekly to monthly intervals.

\subsubsection{CHIRPS Rainfall Data}
Climate Hazards group InfraRed Precipitation with Station (CHIRPS) data provides quasi-global satellite-derived rainfall estimates at sub-national resolution. We extracted monthly rainfall totals aggregated at the sub-national administrative level (adm\_level = 1).

\subsubsection{World Bank Commodity Prices}
Monthly global commodity price indices were obtained from the World Bank's Commodity Markets database. We extracted crude oil average prices (as a transport cost proxy) and global maize prices (as an international supply-demand indicator).

\subsection{Data Preprocessing}

\subsubsection{Temporal Alignment}
All three datasets operate at different temporal frequencies: WFP prices are irregular (weekly/monthly), CHIRPS provides dekadal data, and World Bank indices are strictly monthly. We standardised all data to month-start (MS) frequency using Pandas' \texttt{resample('MS').mean()} function.

\subsubsection{Geographic Alignment}
The WFP dataset uses text-based district names (e.g., ``Kampala'', ``Gulu'') while CHIRPS uses ISO P-codes (``UG3'', ``UG1'', ``UG2'' for Kampala, Gulu, and Mbale regions respectively). We constructed a manual mapping dictionary to align spatial references.

\subsubsection{Missing Value Treatment}
Missing price values were filled using linear interpolation with a maximum gap of 3 consecutive months. Remaining covariates (rainfall, World Bank indices) were handled using forward-fill followed by backward-fill to avoid introducing future information.

\subsubsection{Commodity Standardisation}
In the WFP dataset, ``Maize (white)'' was mapped to ``Maize'' to ensure consistent commodity naming. Prices were aggregated to monthly mean values per unique (date, district, commodity) combination.

\subsection{Feature Engineering}
We engineered five categories of predictive features:

\begin{enumerate}
    \item \textbf{Price Lag Features:} $\text{price\_lag}_k(t) = \text{price}(t - k)$ for $k \in \{1, 2, 4, 8\}$, capturing short- and medium-term price momentum.
    
    \item \textbf{Rolling Statistics:} Rolling mean and standard deviation over 4-month and 8-month windows:
    \begin{equation}
        \text{roll\_mean}_w(t) = \frac{1}{w}\sum_{i=0}^{w-1} \text{price}(t-i)
    \end{equation}
    
    \item \textbf{Seasonal Encoding:} Sinusoidal transformation of the month variable:
    \begin{equation}
        \text{month\_sin}(t) = \sin\left(\frac{2\pi \cdot m(t)}{12}\right), \quad \text{month\_cos}(t) = \cos\left(\frac{2\pi \cdot m(t)}{12}\right)
    \end{equation}
    This preserves the cyclical continuity between December and January, unlike one-hot encoding. Binary indicators for Uganda's two harvest seasons (Season A: June--July; Season B: December--January) were also included.
    
    \item \textbf{Rainfall Lag Features:} $\text{rain\_lag}_k(t) = \text{rainfall}(t - k)$ for $k \in \{1, 4, 8\}$, reflecting the delayed impact of rainfall on crop supply and market prices.
    
    \item \textbf{Global Covariates:} World Bank crude oil and maize price indices included directly.
\end{enumerate}

After feature engineering, rows with NaN values (from lagging operations) were dropped. The final datasets contained approximately 150--180 observations per crop-district pair.

\subsection{Train/Test Split}
A strict chronological 80/20 split was applied: the first 80\% of observations (by date) formed the training set, and the final 20\% formed the test set. No shuffling was applied to preserve temporal ordering and prevent data leakage.

\subsection{Model Implementations}

\subsubsection{Model 1: ARIMA (SARIMAX)}
We implemented ARIMA using the \texttt{statsmodels.tsa.statespace.sarimax.SARIMAX} class. The optimal order $(p, d, q)$ was selected via exhaustive AIC grid search over $p \in \{0, 1, 2, 3\}$, $d \in \{0, 1\}$, $q \in \{0, 1, 2, 3\}$, yielding 32 candidate models per series. The order with the lowest Akaike Information Criterion (AIC) was selected:
\begin{equation}
    \text{AIC} = 2k - 2\ln(\hat{L})
\end{equation}
where $k$ is the number of parameters and $\hat{L}$ is the maximised likelihood. Forecasting used the \texttt{forecast()} method with recursive multi-step prediction.

\subsubsection{Model 2: Facebook Prophet}
Prophet was implemented via the \texttt{prophet} library. Yearly seasonality was enabled with Fourier order $N = 10$. Weekly and daily seasonality were disabled (data is monthly). Rainfall was added as an external regressor via \texttt{add\_regressor('rainfall')}. The model decomposes the time series as:
\begin{equation}
    y(t) = g(t) + s(t) + \beta_{\text{rain}} \cdot \text{rainfall}(t) + \varepsilon_t
\end{equation}
where $g(t)$ is the piecewise linear trend, $s(t)$ is the Fourier seasonal component, and $\beta_{\text{rain}}$ is the learned rainfall coefficient.

\subsubsection{Model 3: LSTM Neural Network}
The LSTM was implemented using TensorFlow/Keras with the following architecture:

\begin{table}[H]
\centering
\caption{LSTM Network Architecture}
\begin{tabular}{lll}
\toprule
\textbf{Layer} & \textbf{Configuration} & \textbf{Output Shape} \\
\midrule
LSTM Layer 1 & 64 units, return\_sequences=True & (batch, 6, 64) \\
Dropout 1 & rate = 0.2 & (batch, 6, 64) \\
LSTM Layer 2 & 32 units, return\_sequences=False & (batch, 32) \\
Dropout 2 & rate = 0.2 & (batch, 32) \\
Dense Output & 1 unit, linear activation & (batch, 1) \\
\bottomrule
\end{tabular}
\end{table}

\textbf{Input Features:} The LSTM receives 9 features: price, rainfall, wb\_crude\_oil, wb\_maize, price\_lag\_1, price\_lag\_2, month\_sin, month\_cos, and rain\_lag\_1.

\textbf{Sliding Window:} Input sequences of length 6 months (window = 6) are constructed. The model examines 6 consecutive monthly observations to predict the price at month 7.

\textbf{Data Normalisation:} All features were scaled to $[0, 1]$ using MinMaxScaler:
\begin{equation}
    x_{\text{scaled}} = \frac{x - x_{\min}}{x_{\max} - x_{\min}}
\end{equation}
Predictions were inverse-transformed to recover original UGX values.

\textbf{Training Configuration:} Adam optimiser (learning rate = 0.001), MSE loss function, batch size = 16, maximum 100 epochs with EarlyStopping (patience = 10, monitor = loss, restore\_best\_weights = True). Random seed set to 42 for reproducibility.

\textbf{Multi-step Forecasting:} For the Streamlit dashboard, the LSTM uses a recursive strategy: after predicting month $t+1$, the prediction is fed back as input, lag features are updated, and the model predicts $t+2$, continuing for the specified horizon.

\subsection{Evaluation Metrics}
Three complementary metrics were used:

\begin{equation}
    \text{MAE} = \frac{1}{n}\sum_{i=1}^{n}|y_i - \hat{y}_i|
\end{equation}
\begin{equation}
    \text{RMSE} = \sqrt{\frac{1}{n}\sum_{i=1}^{n}(y_i - \hat{y}_i)^2}
\end{equation}
\begin{equation}
    \text{MAPE} = \frac{100\%}{n}\sum_{i=1}^{n}\left|\frac{y_i - \hat{y}_i}{y_i}\right|
\end{equation}

MAE provides interpretable error in UGX. RMSE penalises large errors disproportionately. MAPE enables scale-independent cross-crop comparison.

% ============================================================
% 4. RESULTS
% ============================================================
\section{Results}

\subsection{Exploratory Data Analysis}
Before modelling, we conducted comprehensive EDA including: price trend analysis across all 9 crop-district pairs, rainfall pattern visualisation by region, World Bank commodity trend analysis, price distribution box plots, correlation heatmaps, seasonal decomposition, ACF/PACF analysis, and price vs. rainfall scatter plots. Key findings:

\begin{itemize}
    \item Beans exhibited the highest price volatility across all districts (coefficient of variation > 0.4).
    \item Kampala prices were consistently highest, reflecting urban market premiums and transport costs.
    \item Gulu showed distinct post-conflict price patterns with structural breaks around 2016--2017.
    \item Rainfall showed strong annual cyclicality with two clear wet seasons corresponding to Uganda's bimodal pattern.
    \item ACF analysis confirmed significant autocorrelation at lags 1, 2, and 12, supporting our lag feature selection.
\end{itemize}

\subsection{ARIMA Results}
The AIC grid search identified optimal orders for each crop-district pair. ARIMA performed best in districts with more stable, trend-driven price series (e.g., Mbale Maize: MAPE = 12.25\%) and worst for highly volatile series (e.g., Kampala Beans: MAPE = 48.35\%).

\subsection{Prophet Results}
Prophet showed mixed performance. It achieved competitive results for some pairs (e.g., Gulu Cassava flour: MAPE = 12.34\%) where strong annual seasonality was present, but performed poorly when automatic changepoint detection identified spurious trend shifts (e.g., Mbale Beans: MAPE = 52.50\%).

\subsection{LSTM Results}
The LSTM consistently achieved the best or second-best MAPE across all 9 crop-district pairs. Its strongest performance was on series with complex nonlinear dynamics that benefit from multivariate inputs.

\subsection{Comparative Analysis}

\begin{table}[H]
\centering
\caption{Complete Model Comparison: All Crop-District Pairs}
\label{tab:full_results}
\small
\begin{tabular}{llcccc}
\toprule
\textbf{District} & \textbf{Crop} & \textbf{Model} & \textbf{MAE} & \textbf{RMSE} & \textbf{MAPE (\%)} \\
\midrule
Kampala & Maize & ARIMA & 855.93 & 1022.14 & 40.66 \\
Kampala & Maize & Prophet & 768.55 & 917.74 & 37.13 \\
Kampala & Maize & LSTM & 654.56 & 789.20 & 31.33 \\
\midrule
Kampala & Beans & ARIMA & 2093.33 & 2231.08 & 48.35 \\
Kampala & Beans & Prophet & 1505.29 & 1631.09 & 35.59 \\
Kampala & Beans & LSTM & 780.68 & 972.91 & 18.15 \\
\midrule
Kampala & Cassava flour & ARIMA & 739.81 & 874.99 & 29.40 \\
Kampala & Cassava flour & Prophet & 714.82 & 846.10 & 28.40 \\
Kampala & Cassava flour & LSTM & 450.67 & 516.93 & 18.84 \\
\midrule
Gulu & Maize & ARIMA & 188.49 & 230.19 & 13.49 \\
Gulu & Maize & Prophet & 340.74 & 382.42 & 26.39 \\
Gulu & Maize & LSTM & 162.92 & 203.26 & 11.59 \\
\midrule
Gulu & Beans & ARIMA & 849.33 & 1030.07 & 18.60 \\
Gulu & Beans & Prophet & 880.45 & 1059.37 & 16.18 \\
Gulu & Beans & LSTM & 606.74 & 725.40 & 12.22 \\
\midrule
Gulu & Cassava flour & ARIMA & 1033.67 & 1303.28 & 22.67 \\
Gulu & Cassava flour & Prophet & 561.22 & 844.64 & 12.34 \\
Gulu & Cassava flour & LSTM & 878.20 & 1152.36 & 20.01 \\
\midrule
Mbale & Maize & ARIMA & 169.37 & 191.66 & 12.25 \\
Mbale & Maize & Prophet & 223.11 & 269.37 & 18.37 \\
Mbale & Maize & LSTM & 193.93 & 226.59 & 14.22 \\
\midrule
Mbale & Beans & ARIMA & 409.95 & 518.76 & 11.52 \\
Mbale & Beans & Prophet & 1936.23 & 2002.98 & 52.50 \\
Mbale & Beans & LSTM & 443.49 & 484.38 & 11.78 \\
\midrule
Mbale & Cassava flour & ARIMA & 477.18 & 538.77 & 29.08 \\
Mbale & Cassava flour & Prophet & 641.05 & 774.99 & 40.07 \\
Mbale & Cassava flour & LSTM & 356.97 & 386.56 & 20.75 \\
\bottomrule
\end{tabular}
\end{table}

\begin{table}[H]
\centering
\caption{Average Metrics by Model (Across All 9 Crop-District Pairs)}
\label{tab:avg_results}
\begin{tabular}{lccc}
\toprule
\textbf{Model} & \textbf{Avg MAE (UGX)} & \textbf{Avg RMSE} & \textbf{Avg MAPE (\%)} \\
\midrule
ARIMA & 757.45 & 882.33 & 25.11 \\
Prophet & 841.27 & 969.86 & 29.66 \\
\textbf{LSTM} & \textbf{503.13} & \textbf{606.40} & \textbf{17.65} \\
\bottomrule
\end{tabular}
\end{table}

The LSTM achieved the lowest average MAE (503.13 UGX), RMSE (606.40), and MAPE (17.65\%), outperforming ARIMA by 33\% and Prophet by 40\% on the MAE metric.

% ============================================================
% 5. DISCUSSION
% ============================================================
\section{Discussion}

\subsection{Key Findings}

\subsubsection{LSTM Superiority}
The LSTM's dominance stems from three capabilities: (1) nonlinear modelling capacity -- agricultural price dynamics involve threshold effects and regime changes that linear ARIMA cannot capture; (2) multivariate input processing -- the LSTM jointly considers 9 features including climate and global commodity data, while ARIMA operates on univariate price data; (3) long-range memory -- the gating mechanism preserves information across the 6-month sliding window, capturing seasonal interactions that Prophet's fixed Fourier basis may miss.

\subsubsection{Prophet Underperformance}
Prophet's higher MAPE (29.66\% vs ARIMA's 25.11\%) was unexpected. Three factors explain this: (a) automatic changepoint detection identified spurious trend breaks in our relatively short series, particularly for Mbale Beans where Prophet's MAPE reached 52.50\%; (b) Prophet's Fourier seasonality assumes fixed seasonal timing, while Uganda's harvest periods shift by 2--3 weeks annually depending on rainfall onset; (c) Prophet's rainfall regressor assumes a constant linear coefficient, while the true price-rainfall relationship varies by crop growth stage.

\subsubsection{Geographic Variation}
Models performed best in Gulu and Mbale (lower MAPEs) compared to Kampala. This is counterintuitive -- one might expect the largest market to be most predictable. However, Kampala's role as a national hub means its prices are influenced by nationwide supply dynamics, speculative trading, and import/export flows, creating more complex nonlinear patterns. Gulu and Mbale, being more locally-driven markets, exhibit more regular seasonal patterns amenable to time-series modelling.

\subsection{Limitations}

\begin{enumerate}
    \item \textbf{Data Scarcity:} With approximately 150--180 observations per series, our LSTM operates at the lower bound of viable training data for neural networks. More historical data would likely improve performance.
    
    \item \textbf{Univariate ARIMA:} Our ARIMA implementation used only price data. SARIMAX with exogenous variables (rainfall, oil prices) could improve ARIMA's competitiveness.
    
    \item \textbf{Single Train/Test Split:} We used one chronological split rather than time-series cross-validation with expanding windows, which would provide more robust performance estimates.
    
    \item \textbf{Static Features:} Climate covariates were used as-is without incorporating forecasted weather, which would be necessary for a production deployment.
    
    \item \textbf{No Confidence Intervals:} Our point forecasts lack uncertainty quantification, which is critical for risk-sensitive decision making.
\end{enumerate}

\subsection{Implications for Food Security}
A forecasting system with 17\% MAPE could meaningfully support food security decisions:
\begin{itemize}
    \item \textbf{For farmers:} Knowing that maize prices are predicted to rise 15\% in 6 weeks provides actionable intelligence for storage decisions.
    \item \textbf{For NGOs:} Early warning of price spikes enables pre-positioning of food aid and cash transfer programmes.
    \item \textbf{For policymakers:} District-level price forecasts inform decisions about strategic grain reserves and import tariff adjustments.
\end{itemize}

% ============================================================
% 6. DEPLOYMENT
% ============================================================
\section{System Deployment}

\subsection{Streamlit Dashboard}
The best-performing LSTM model was deployed in an interactive web dashboard built with Streamlit. The dashboard provides:

\begin{itemize}
    \item \textbf{User Interface:} Sidebar with dropdown selectors for crop, district, and forecast horizon (4 or 8 weeks ahead).
    \item \textbf{Real-time Training:} The LSTM model trains on-the-fly when parameters are selected, using Streamlit's caching mechanism for efficiency.
    \item \textbf{Visualisation:} Interactive Plotly charts displaying historical data, test-set evaluation predictions, and future forecast trajectory.
    \item \textbf{Performance Metrics:} Displayed MAE, RMSE, and MAPE for the current crop-district selection.
    \item \textbf{Model Comparison Panel:} Side-by-side comparison of ARIMA, Prophet, and LSTM metrics with bar chart visualisations.
\end{itemize}

\subsection{Technical Implementation}
The dashboard loads and preprocesses data from CSV/Excel files, applies the same feature engineering pipeline as the notebook, trains the LSTM with identical hyperparameters, and generates recursive multi-step forecasts for the selected horizon.

% ============================================================
% 7. CONCLUSIONS AND FUTURE WORK
% ============================================================
\section{Conclusions and Future Work}

\subsection{Conclusions}
This project successfully built and compared three fundamentally different forecasting paradigms for Ugandan staple crop prices. The LSTM neural network achieved the best overall performance with a 17.65\% average MAPE, demonstrating that multivariate deep learning approaches incorporating climate data can meaningfully outperform classical statistical methods in data-scarce agricultural contexts. The complete pipeline -- from data collection through model training to interactive deployment -- demonstrates the full lifecycle of an applied machine learning project.

\subsection{Future Work}
\begin{enumerate}
    \item \textbf{Expanded Coverage:} Extend to additional districts and crops, potentially covering all 135 Ugandan districts.
    \item \textbf{Transformer Architectures:} Implement attention-based models (e.g., Temporal Fusion Transformer) for potentially improved long-range forecasting.
    \item \textbf{Ensemble Methods:} Combine ARIMA, Prophet, and LSTM predictions with learned weights.
    \item \textbf{Additional Covariates:} Incorporate conflict data, COVID-19 mobility indices, and soil moisture indicators.
    \item \textbf{Probabilistic Forecasting:} Generate prediction intervals rather than point forecasts.
    \item \textbf{Cloud Deployment:} Deploy on AWS/GCP for public stakeholder access.
    \item \textbf{SMS Integration:} Build a USSD/SMS interface for farmer access via feature phones.
\end{enumerate}

% ============================================================
% REFERENCES
% ============================================================
\newpage
\begin{thebibliography}{99}

\bibitem{box2015}
Box, G.E.P., Jenkins, G.M., Reinsel, G.C. and Ljung, G.M. (2015). \textit{Time Series Analysis: Forecasting and Control}. 5th ed. Hoboken, NJ: John Wiley \& Sons.

\bibitem{brown2009}
Brown, M.E. (2009). ``Markets, Climate Change, and Food Security in West Africa.'' \textit{Environmental Science \& Technology}, 43(21), pp.8016--8020.

\bibitem{hochreiter1997}
Hochreiter, S. and Schmidhuber, J. (1997). ``Long Short-Term Memory.'' \textit{Neural Computation}, 9(8), pp.1735--1780.

\bibitem{sabu2020}
Sabu, K.M. and Kumar, T.K.M. (2020). ``Predictive Analytics in Agriculture: Forecasting Prices of Arecanuts Using ARIMA and LSTM.'' \textit{Procedia Computer Science}, 171, pp.699--708.

\bibitem{taylor2018}
Taylor, S.J. and Letham, B. (2018). ``Forecasting at Scale.'' \textit{The American Statistician}, 72(1), pp.37--45.

\bibitem{worldbank2023}
World Bank (2023). \textit{Uganda Economic Update: Agricultural Sector Review}. Washington, DC: World Bank Group.

\bibitem{xiong2015}
Xiong, R., Nichols, E.P. and Shen, Y. (2015). ``Deep Learning for Stock Market Prediction from Financial News Articles.'' \textit{IEEE International Conference on Big Data}, pp.2015--2021.

\bibitem{zhang2003}
Zhang, G.P. (2003). ``Time Series Forecasting Using a Hybrid ARIMA and Neural Network Model.'' \textit{Neurocomputing}, 50, pp.159--175.

\end{thebibliography}

\end{document}
